\documentclass[a4paper]{amsart} 
\usepackage{amssymb} 
\usepackage{hyperref}
\pagestyle{empty}
\date{} 

%% IGNORE THE NEXT 4 LINES
\makeatletter % 
\def\labelenumi{\theenumi.}
\def\theenumi{\@arabic\c@enumi}
\makeatother
%% STOP IGNORING NOW

\begin{document}

\title{APPLICATIONS OF COMPANION FORMS TO EIGENVALUE BOUNDS AND SCALAR POLYNOMIALS}
\author{Aaron Melman} %Presenting author only
\address{Santa Clara University} % Affiliation of presenting author
\email{amelman@scu.edu} % Email address of presenting author only

\maketitle

% The abstract  ***no more than one page***

We present a simple way to derive companion forms of matrix polynomials, which are lower order matrix polynomials with the same eigenvalues 
(so-called "$\ell$-ifications") as the given matrix polynomial, and show how they can be used to produce eigenvalue bounds.
These bounds, which also include nonstandard directional ones, can be substantially less computationally demanding when using higher degree companion 
forms, as opposed to classical linearizations (companion forms of degree one).

As an application to scalar polynomials, we show how companion forms provide a convenient way for the derivation of a polynomial whose (unknown) zeros
are powers of those of a given polynomial.




\end{document}


% Please leave the next bit alone  
%%% Local Variables: 
%%% mode: latex
%%% TeX-master: "../ILAS-2022"
%%% End:
